% Textos fijos, espa\~nol. Para un juego, sustituye alguno en override-es.tex (\renewcommand).
\usepackage[spanish]{babel}
\renewcommand{\portWork}{\portWorkES}
\newcommand{\tManual}{Manual de usuario}
\newcommand{\tUserManual}{MANUAL DE USUARIO}
\newcommand{\tForV}{para V9990}
\newcommand{\tForTheV}{para el Yamaha V9990}
\newcommand{\tCoverLine}{El juego de MSX en el V9990}
\newcommand{\tOriginalBy}{Un juego original de}
\newcommand{\tPortBy}{Versi\'on para V9990 de}
\newcommand{\tHomage}{Esta es una versi\'on no oficial hecha para el estudio, la conservaci\'on y como homenaje.}
\newcommand{\tWhatYouNeed}{Qu\'e necesitas}
\newcommand{\tChipCaption}{El procesador de v\'ideo Yamaha V9990, el chip con el que dibuja esta versi\'on del juego.}
\newcommand{\tItem}{Elemento}
\newcommand{\tRequirement}{Requisito}
\newcommand{\tComputer}{Ordenador}
\newcommand{\tComputerText}{\textbf{Cualquier MSX1} con dos ranuras de cartucho, una para el juego y otra para el V9990. Probado en un Sony HB-10P (MSX1) y en un MSX turbo R; otras generaciones MSX deber\'ian funcionar pero no se han probado.}
\newcommand{\tMemory}{Memoria}
\newcommand{\tMemoryText}{\textbf{\ramNeeded{} de RAM} en la memoria principal (las cuatro p\'aginas).}
\newcommand{\tVideo}{V\'ideo}
\newcommand{\tVideoText}{Un \textbf{cartucho V9990} (v\'ideo Yamaha V9990). Es obligatorio: todo el dibujo lo hace el V9990. El chip de v\'ideo del propio ordenador se deja en blanco a prop\'osito.}
\newcommand{\tSoundRow}{Sonido}
\newcommand{\tSoundText}{Siempre se usa el PSG del MSX. Un \textbf{SCC} (en el cartucho del juego o en otra ranura) es opcional.}
\newcommand{\tGameRow}{Juego}
\newcommand{\tGameText}{El \textbf{cartucho de \gameTitle{} para V9990}: una \romKind{} (el fichero \texttt{\romFile}, grabado en un cartucho flash o equivalente).}
\newcommand{\tMonitor}{Monitor}
\newcommand{\tMonitorText}{Un monitor o televisor conectado a la \textbf{salida de v\'ideo del V9990}, no a la salida de v\'ideo del propio MSX.}
\newcommand{\tControlsRow}{Controles}
\newcommand{\tControlsText}{Teclado, cursores o un joystick en el puerto 1.}
\newcommand{\tNoTape}{\textbf{Sin cinta, sin disco, sin contrase\~na.} Todo el juego est\'a en el cartucho.}
\newcommand{\tSettingUp}{Puesta en marcha}
\newcommand{\tSetOff}{\textbf{Apaga el MSX.} No insertes ni retires nunca un cartucho con el equipo encendido.}
\newcommand{\tSetV}{Inserta el \textbf{cartucho V9990} en una de las ranuras.}
\newcommand{\tSetGame}{Inserta el \textbf{cartucho de \gameTitle{}} en la otra ranura. El orden no importa.}
\newcommand{\tSetMonitor}{Conecta el monitor o televisor a la \textbf{salida de v\'ideo del V9990}.}
\newcommand{\tSetJoy}{Conecta un joystick al puerto 1 si quieres usarlo.}
\newcommand{\tSetOn}{Enciende el MSX. Un momento la pantalla del propio MSX muestra \texttt{V9990 found !}; despu\'es la imagen aparece en el monitor del V9990. Si dice \texttt{V9990 NOT found!!!} falta el V9990 o no est\'a bien insertado, y el juego se detiene ah\'i.}
\newcommand{\tBlackScreen}{\textbf{Si la pantalla se queda en negro} comprueba, por este orden: que el monitor est\'a en la salida del V9990; que el V9990 y el juego est\'an bien insertados; que el ordenador tiene \ramNeeded{} de RAM; que la ROM del juego es la versi\'on de \romKind. En un emulador como openMSX, activa la extensi\'on V9990 y carga la ROM con el mapper correcto.}
\newcommand{\tPowerOnTitle}{Del encendido al juego}
\newcommand{\tPowerOnCaption}{Izquierda: la portada. Derecha: la pantalla de t\'itulo.}
\newcommand{\tControlsTitle}{Controles}
\newcommand{\tCreditsTitle}{5. Cr\'editos y derechos}
\newcommand{\tIsAGameBy}{es un juego de}
\newcommand{\tRightsText}{El juego, sus gr\'aficos, m\'usica y nombre pertenecen a sus titulares. Esta es una versi\'on no oficial hecha para el estudio, la conservaci\'on y como homenaje; no tiene relaci\'on con ellos.}
\newcommand{\tVersionWord}{Versi\'on V9990}
\newcommand{\tMsxTrademark}{MSX es una marca de MSX Licensing Corporation. El logotipo MSX se muestra solo para indicar en qu\'e ordenadores funciona el juego}
\newcommand{\tChipNote}{V9990 es el nombre del chip de v\'ideo de Yamaha; la foto del chip es de Yaca2671, CC BY-SA 3.0.}
\newcommand{\tSourceNote}{C\'odigo fuente disponible en GitHub. \textbf{\textexclamdown Mejora el juego si te apetece!} (\textit{Please improve the game if you feel so!})}
\newcommand{\tAppendix}{Ap\'endice. Etiqueta del cartucho}
\newcommand{\tStickerText}{Una etiqueta para el cartucho, lista para imprimir. Al 100\,\% mide unos 16\,cm de ancho; esc\'alala en tu programa de impresi\'on al \'area de etiqueta de tu cartucho (la imagen tiene proporci\'on 3:2).}
\newcommand{\tAction}{Acci\'on}
\newcommand{\tKeyboard}{Teclado}
\newcommand{\tJoystick}{Joystick}
\newcommand{\tNo}{N.\textordmasculine}
\newcommand{\tElement}{Elemento}
\newcommand{\tTellsYou}{Qu\'e te indica}
\newcommand{\tTips}{Consejos.}
